\section{Numerical coding and the induction discipline}\label{sec:coding}

This appendix fixes enough representation detail to make the route from
the mathematical arguments to first-order iPRA unambiguous. The particular
codes can be changed by proved PR translations. Their logical role cannot
be omitted.

\subsection{Total numerical operations and validity guards}

Use the signed coding
\[
 \operatorname{enc}(z)=
 \begin{cases}2z,&z\ge0,\\2|z|-1,&z<0.\end{cases}
\]
Decoding, sign, absolute value, comparison and the integer ring operations
are PR. Their laws reduce to natural addition, multiplication
and truncated subtraction. A rational is a pair $(z,d)$ with $d>0$;
canonical reduction by the gcd is convenient but unnecessary for the
arithmetic estimates. Equality and order use cross multiplication with
positive denominators. On a malformed rational code, use a fixed default
value and keep the validity flag explicit. Division in the ODE algorithm
always has a separately proved positive denominator.

For a concrete pairing, one can use
$\langle a,b\rangle=2^a(2b+1)-1$. Bounded trial division and a bounded
search recover its components. The decoder is total. A finite sequence
can be coded by a length together with nested pairs. Access to entry $i$
is a recursion bounded by the length, returning a default past the end.
Alternatively, a table whose entries are bounded by $C$ can be coded in
base $C+1$. Then a length-$r$ table has code below $(C+1)^r$.
Both encodings and the usual append, map, fold and bounded search
operations are PR.

Canonical polynomial codes remove trailing zero coefficients. In the
split-branch checker, the remainder must be nonzero, and the chosen
difference is recomputed from the two selected input polynomials.
Comparison of normalized lengths expresses the strict degree gap even
for constant remainders. The zero polynomial is not assigned a genuine
degree and then accidentally admitted as a certificate.

\subsection{Why variable size does not introduce a universal evaluator}

The interpreter for a finite integer or rational polynomial table is one
fixed PR program. It performs bounded folds over coefficients, monomials
and coordinates, using numerical exponentiation. It does not evaluate
arbitrary PR program descriptions. Iterating this fixed update a variable
number of times is precisely the primitive recursion operation allowed
in the language.

For example, let $s_{j+1}=\operatorname{Step}(s_j,c)$ where $c$ is a
polynomial-code parameter. The function $(j,c,s_0)\mapsto s_j$ is PR
because \code{Step} is one fixed PR function, even if its outputs become
very large. Its evaluation time need not be elementary or practical.
Nothing here asserts that the evaluator of all PR syntax is itself PR.

The same distinction applies to Lean's function-valued notation
\code{Fin d -> Rat}. This is a convenient finite coordinate presentation
in the metatheory. The object theory uses its explicit finite sequence
code, and the correspondence is proved by entrywise decoding lemmas.
Arbitrary functions $\N\to\Q$ are never inserted into the numerical
state as if they were finite data.

\subsection{Bounded formulas and finite induction}

For a fixed PR Boolean test $b(\vec x,i)$, define
\[
 A(\vec x,0)=1,\qquad
 A(\vec x,j+1)=A(\vec x,j)\mathbin{\wedge_b}b(\vec x,j),
\]
where $\wedge_b$ is conjunction on $0,1$. A finite induction
proves that $A(\vec x,j)=1$ expresses success at every index $i<j$.
For the forward direction, fix an arbitrary challenge $i$ and induct on
$j$ with the quantifier-free invariant
$(i<j\land A(\vec x,j)=1)\to b(\vec x,i)=1$.
For the reverse direction, fix a terminal bound $N$ and keep
$\forall i<N\,b(\vec x,i)=1$ as an assumption independent of the
induction variable. Induct on $j$ with
$j\le N\to A(\vec x,j)=1$. Universal elimination from the fixed
assumption supplies the successor test. Thus even the proof of bounded
reflection does not presuppose induction on a quantified formula.
The corresponding disjunctive fold and first-success register implement
bounded existence and witness selection. With atomic equality and order
tests, structural induction on a fixed bounded formula proves its
equivalence to a Boolean equation. This is the reflection interface used
throughout; it is not reflection for all theorems of iPRA.

If a bounded loop has guard $j\le N$, its invariant can have the form
\[
 (j\le N\longrightarrow B(s_j,j,\vec x)=1).
\]
Here order is a PR test. A stopped transition freezes the state outside
the active range. The induction formula is quantifier-free. In particular,
the assertion that the first $j$ entries all satisfy a bound must be
expressed by its Boolean prefix fold before invoking \code{PRA.qfInd}.

The principal invariants in the new applications are:
\begin{longtable}{@{}>{\raggedright\arraybackslash}p{.22\linewidth}>{\raggedright\arraybackslash}p{.42\linewidth}>{\raggedright\arraybackslash}p{.30\linewidth}@{}}
\toprule
Loop & Numerical invariant & Necessary guard or auxiliary data \\
\midrule\endhead
Polynomial evaluation & Horner prefix agrees with the corresponding finite
sum & Encoded coefficient length and prefix index \\
Polynomial bound & Absolute partial sum is bounded by the sum of absolute
terms & Nonnegative power base; leading coefficient nonzero when used \\
Euler trajectory & $\norm{x_j-x_0}\le jhM$ & $j\le N$,
$hN=T$, and the box bound \\
Error propagation & $e_j\le r^j(e_0+\sum_{i<j}\delta_i)$ & $r\ge1$
and nonnegative rational errors \\
Conservation & Accumulated discrepancy is at most the sum
of quadratic step remainders & Coefficientwise Lie-derivative identity \\
Split factors & Each selected nonzero integer factor has absolute value
at most the product & Full product nonzero; zero-product case separate \\
Rectangle scan & Output contains exactly the successful tests already
visited & Explicit bounds $H,K$ and scan position \\
\bottomrule
\end{longtable}

\subsection{Accuracy challenges and rational gap elimination}

An equality of names is an assertion about all requested precisions. The
correct proof order is to fix an arbitrary precision $k$, compute a finite
schedule from $k$, prove the corresponding rational inequality by the
preceding induction discipline, and then universally introduce $k$.
One does not induct on equality of represented values.

Removing a vanishing numerical error is elementary. Suppose
rational $a,b,C$ with $C\ge0$ satisfy $a\le b+C2^{-n}$ for all $n$.
If $a>b$, write the positive rational gap $a-b=u/v$ with $u,v\ge1$.
Choose an integer $J\ge Cv/u$ and take $n=J+1$. The finite inequality
$2^{J+1}>J$ makes $C2^{-n}<u/v$ when $C>0$, a contradiction.
For $C=0$ the conclusion is immediate. The choice of $J$ is an explicit
rational ceiling, and the exponent bound is PR. This removes a rational
slack; it does not decide the sign of an arbitrary represented real.

Finally, every theorem described as uniform has one finite numerical
algorithm and one arithmetic proof schema for that algorithm. A collection
of unrelated standard-numeral proofs is not substituted for a single
variable-parameter proof.
